%%% Local Variables:
%%% mode: latex
%%% TeX-master: "../main"
%%% End:

\chapter[精细平衡类步态专家的独立训练]{精细平衡类步态专家的独立训练}
\label{cha:fine-experts}

本章对应"分训"阶段的第二类专家——精细平衡类步态专家，阐述抬腿行走专家、抬腿上台阶专家与单腿平衡专家三个专家网络的独立训练。三者均要求在大范围抬腿/重心转移的同时保持姿态稳定，涉及"高动态动作"与"精细平衡"的矛盾，是精细控制与约束强化学习的典型应用。

\section{精细平衡类专家库概述}

精细平衡类专家库包含三个专家：

（1）\textbf{抬腿行走专家}：面向粗糙与木桩地形，通过周期抬腿实现"点足"行走，为越障与精细化操控提供步态基础，采用PPO训练；

（2）\textbf{抬腿上台阶专家}：面向低台阶地形，实现稳定抬腿跨越台阶，采用N3PO约束强化学习训练；

（3）\textbf{单腿平衡专家}：面向独木桥与狭窄地形，实现单侧轮支撑平衡行走，采用PPO训练。

该类任务的共同难点在于约束苛刻且动作与稳定目标冲突剧烈，若不显式处理约束，策略极易在"抬腿高度"与"平衡保持"之间顾此失彼，或陷入剧烈晃动的局部最优。

\section{抬腿行走专家}

\subsection{抬腿机制与任务设计}

抬腿行走通过周期性地抬起单侧腿、以另一侧轮/腿支撑实现"点足"行走。其观测空间在基础观测上追加腿部抬升高度、相位计时器与轮地接触标志等维度。任务定义与滚动行走类似，但以"抬腿高度跟踪"与"支撑稳定性"作为核心奖励项。

\subsection{粗糙与木桩地形的点足行走}

在粗糙地形上，滚动行走轮端颠簸明显、越障能力有限，而抬腿点足行走能够主动选择支撑点、减小轮端无效接触。在木桩地形上，机器人需在有限离散支撑点之间规划落脚位置。因此，本文将该专家设计为调度器在粗糙与木桩地形的首选选项——这也是第3章表~\ref{tab:terrain-gait}中"粗糙/碎石地形→抬腿点足行走"对应关系的实现基础。

\section{抬腿上台阶专家}

\subsection{任务定义与约束分析}

抬腿上台阶要求机器人在保持稳定站立与前进后退转向能力的前提下，将单侧腿抬起并跨越台阶边缘，完成上台阶与下台阶动作。任务的关键约束包括：关节力矩饱和限制、轮端净空要求、机身姿态范围限制以及落地冲击限制。

\subsection{基于N3PO的约束强化学习}

本章将上台阶任务建模为约束马尔可夫决策过程（CMDP）。在总奖励中保留基础能力奖励项的同时，将关键安全指标建模为代价函数并纳入N3PO的约束优化框架
\begin{equation}
  \begin{aligned}
    \max_{\pi}\;&\mathbb{E}\left[\sum_t \gamma^t\left(r_{task} + w\,r_{base}\right)\right],\\
    \text{s.t.}\;&\mathbb{E}\left[\sum_t \gamma^t\, c_{sat}\right] \le d_1,\quad
    \mathbb{E}\left[\sum_t \gamma^t\, c_{tilt}\right] \le d_2,
  \end{aligned}
  \label{eq:n3po-cmdp}
\end{equation}
式中，$c_{sat}$为关节力矩饱和度代价，$c_{tilt}$为机身倾角越界代价，$d_1$、$d_2$为对应预算。N3PO基于自然策略梯度在每次更新中显式满足约束，保证策略在优化抬腿性能的同时不突破安全边界。

\subsection{与SAC及单强化学习算法的对比}

为验证N3PO的有效性，本章在同一上台阶任务上实现SAC算法与单强化学习算法（纯PPO）作为对比基线。对比指标包括上/下台阶成功率、力矩饱和度、姿态越界率与训练收敛速度，重点考察约束强化学习在"高动态 + 强约束"场景下的优势。

\section{单腿平衡专家}

\subsection{收膝折腿与侧倾机制}

单腿平衡指机器人将重心完全转移至单侧支撑轮，另一侧轮抬离地面，从而获得更窄的等效支撑宽度。为降低支撑难度，本文采用\textbf{收膝折腿}机制：支撑腿收膝下蹲降低质心高度，抬离腿收膝折叠，同时机身向支撑侧侧倾（约$28^\circ$），使质心投影落入支撑轮接触区域内。

与后空翻的"弹道可开环"不同，单腿平衡的保持本质上是倒立摆调节，必须闭环控制。因此，前馈（FF）仅用于折腿与落腿的过渡阶段，横滚平衡等核心调节完全交由强化学习策略闭环实现，其横滚角奖励定义为
\begin{equation}
  r_{roll} = \exp\left(-\lambda\,\left(\frac{up\_y - \cos\theta_{ref}}{\cos\theta_{ref}}\right)^2\right),
  \label{eq:rroll}
\end{equation}
式中，$up\_y$为机身法向在世界系$y$方向的分量，$\theta_{ref}$为参考侧倾角。该奖励在直立时接近0、达到参考侧倾时接近1，引导策略平滑学习侧倾平衡。

\subsection{独木桥行走}

在单腿平衡基础上，本章将步态迁移至单木桥地形：桥面窄、两侧无支撑，机器人需在保持单腿平衡的同时完成前进与转向，对横滚稳定裕度与转向精度要求极高。训练采用热启动方式，将滚动行走专家权重迁移至单腿平衡任务，并通过观测网络输入层扩展（追加FSM状态、计时器与轮地接触历史）适配新任务观测，实现知识迁移与网络扩展。

\section{实验结果与分析}

【待填入：（1）抬腿行走专家在粗糙/木桩地形上的抬腿高度跟踪误差、支撑稳定性、通过率；（2）抬腿上台阶的上/下台阶成功率、力矩饱和度与姿态越界率，N3PO与SAC、单强化学习算法的对比结果；（3）单腿平衡保持时长、横滚角标准差、独木桥通过率与转向精度实验结果，建议以表~\ref{tab:fine-result}汇总对比指标。】

\begin{table}[!htbp]
    \centering
    \bicaption{精细平衡专家对比结果（示意表格，填入真实数据）}{Comparison results of fine balance experts (illustrative, fill in real data)}
    \resizebox{0.92\textwidth}{!}{
        \begin{tabular}{lccc}
            \toprule[1.5pt]
            指标 & N3PO & SAC & 单强化学习 \\
            \midrule[0.75pt]
            上台阶成功率(\%) & 【待填】 & 【待填】 & 【待填】 \\
            力矩饱和度 & 【待填】 & 【待填】 & 【待填】 \\
            单腿平衡时长(s) & 【待填】 & 【待填】 & 【待填】 \\
            独木桥通过率(\%) & 【待填】 & 【待填】 & 【待填】 \\
            \bottomrule[1.5pt]
        \end{tabular}
    }
    \label{tab:fine-result}
\end{table}

\section{本章小结}

本章阐述了精细平衡类步态专家（抬腿行走、抬腿上台阶、单腿平衡）的独立训练方法。抬腿行走专家通过点足抬腿实现粗糙与木桩地形的越障行走；抬腿上台阶专家通过N3PO约束强化学习，在满足力矩与姿态约束的前提下实现稳定上台阶；单腿平衡专家通过收膝折腿与机身侧倾机制，结合热启动迁移实现独木桥行走。三者作为调度器在精细地形场景的专家选项，为统一系统集成提供了精细化操控能力。
